\subsection{自由代数}

定义 2. 设 I 是一个集合；令 M(I)（相应地 Mo(I)，相应地 N \(^{(I)}\) ）表示由 I 导出的自由广群（相应地自由幺半群，相应地自由交换幺半群）。M(I)（相应地 Mo(I)，相应地 N \(^{(I)}\) ）在 A 上的代数称为集合 I 在环 A 上的自由代数（相应地自由结合代数，相应地自由交换结合代数（或作为语言上的滥用，称为自由交换代数））。

我们把 I 在 A 上的自由代数（相应地自由结合代数，相应地自由交换代数）记作 \(\operatorname{Lib}_{\mathbf{A}}(\mathbf{I})\) （相应地 \(\operatorname{Libas}_{\mathbf{A}}(\mathbf{I})\) ，相应地 \(\operatorname{Libasc}_{\mathbf{A}}(\mathbf{I})\) ）。把 I 到 \(\operatorname{M}(\mathbf{I})\) （相应地 \(\operatorname{Mo}(\mathbf{I})\) ，相应地 \(\mathbf{N}^{(\mathbf{I})}\) ）中的典范映射与 \(\operatorname{M}(\mathbf{I})\) （相应地 \(\operatorname{Mo}(\mathbf{I})\) ，相应地 \(\mathbf{N}^{(\mathbf{I})}\) ）到 \(\operatorname{Lib}_{\mathbf{A}}(\mathbf{I})\) （相应地 \(\operatorname{Libas}_{\mathbf{A}}(\mathbf{I})\) ，相应地 \(\operatorname{Libasc}_{\mathbf{A}}(\mathbf{I})\) ）中的典范映射复合，就得到 I 到 \(\operatorname{Lib}_{\mathbf{A}}(\mathbf{I})\) （相应地 \(\operatorname{Libas}_{\mathbf{A}}(\mathbf{I})\) ，相应地 \(\operatorname{Libasc}_{\mathbf{A}}(\mathbf{I})\) ）中的一个典范映射，当 \(A \neq \{0\}\) 时它是单射。我们把元素 \(i \in I\) 在此典范映射下的像记作 \(X_{i}\) ，并称 \(X_{i}\) 为 \(\operatorname{Lib}_{\mathbf{A}}(\mathbf{I})\) （相应地 \(\operatorname{Libas}_{\mathbf{A}}(\mathbf{I})\) ，相应地 \(\operatorname{Libasc}_{\mathbf{A}}(\mathbf{I})\) ）的指标为 i 的未定元。

由于 \(\mathbf{Mo}(\mathbf{I})\) 与 \(\mathbf{N}^{(\mathrm{I})}\) 各有一个单位元， \(\operatorname{Libas}_{\mathbf{A}}(\mathbf{I})\) 与 \(\operatorname{Libasc}_{\mathbf{A}}(\mathbf{I})\)

是含单位元的结合代数，且进一步 \(\operatorname{Libasc}_{\mathbf{A}}(\mathbf{I})\) 是交换的。若 e 是 \(\operatorname{Libas}_{\mathbf{A}}(\mathbf{I})\) （相应地 \(\operatorname{Libasc}_{\mathbf{A}}(\mathbf{I})\) ）的单位元，则映射 \(\alpha \mapsto \alpha e\) 是 A 到 \(\operatorname{Libas}_{\mathbf{A}}(\mathbf{I})\) （相应地 \(\operatorname{Libasc}_{\mathbf{A}}(\mathbf{I})\) ）的中心的一个子环上的同构，该子环与 A 等同（第 1 号）。

命题 7. 设 I 是一个集合，F 是 A 上的一个代数（相应地，含单位元结合代数，相应地，含单位元交换结合代数）。对每个映射 \(f:\mathbf{I}\to\mathbf{F}\) ，存在且仅存在一个同态（相应地，保单位同态） \(\bar{f}\) ，它由 \(\operatorname{Lib}_{\mathbf{A}}(\mathbf{I})\) （相应地 \(\operatorname{Libas}_{\mathbf{A}}(\mathbf{I})\) ，相应地 \(\operatorname{Libasc}_{\mathbf{A}}(\mathbf{I})\) ）到 F 中，使得 \(\bar{f}(\mathbf{X}_{i})=f(i)\) 对所有 \(i\in\mathbf{I}\) 成立。

设 \(F_{m}\) 是通过赋予集合 F 其乘法合成律而得到的广群（相应地，幺半群）。存在且仅存在一个同态（相应地，保单位同态）g，由 M(I)（相应地，Mo(I)，相应地，N \(^{(I)}\) ）到 \(F_{m}\) 中，使得 \(g(i) = f(i)\) 对所有 \(i \in I\) 成立（I，§ 7，第 1、2 与 7 号）；命题 7 随后由第 6 号命题 6 得出。

注记。(1) 我们稍后将定义 \(\operatorname{Libas}_{\mathbf{A}}(\mathbf{I})\) 到自由模 \(\mathbf{A}^{(I)}\) 的张量代数上的一个同构（ \(\S 5\) ，第 5 号），以及 \(\operatorname{Libasc}_{\mathbf{A}}(\mathbf{I})\) 到 \(\mathbf{A}^{(I)}\) 的对称代数上的一个同构（ \(\S 6\) ，第 6 号）。

\begin{enumerate}
\def\labelenumi{(\arabic{enumi})}
\setcounter{enumi}{1}
\item
  设 \(\rho\) 是 A 到交换环 B 中的保单位同态。如已见（§ 2，第 6 号），一个同构 \(\sigma\) 由 \(\rho\) 导出，它把 \(\operatorname{Lib}_{\mathbf{B}}(\mathbf{I})\) （相应地 \(\operatorname{Libas}_{\mathbf{B}}(\mathbf{I})\) ，相应地 \(\operatorname{Libasc}_{\mathbf{B}}(\mathbf{I})\) ）映到代数 \((\operatorname{Lib}_{\mathbf{A}}(\mathbf{I}))_{(\mathbf{B})}\) （相应地 \((\operatorname{Libas}_{\mathbf{A}}(\mathbf{I}))_{(\mathbf{B})}\) ，相应地 \((\operatorname{Libasc}_{\mathbf{A}}(\mathbf{I}))_{(\mathbf{B})}\) ）上，该代数是把标量借助 \(\rho\) 扩张到 B 而得到的；若 \(\mathbf{X}_i^{\mathbf{A}}\) , \(\mathbf{X}_i^{\mathbf{B}}\) 是分别对应于 A 与 B 的指标为 \(i\) 的未定元，则 \(\sigma(\mathbf{X}_i^{\mathbf{B}}) = \mathbf{X}_i^{\mathbf{A}} \otimes 1\) 。
\item
  设 \(J\) 是 \(I\) 的一个子集；我们知道 \(M(J)\) 被等同于广群 \(M(I)\) 的一个稳定子集，因此（第 6 号） \(\operatorname{Lib}_{A}(J)\) 典范地等同于 \(\operatorname{Lib}_{A}(I)\) 的一个子代数，它由诸 \(X_i\) ——满足 \(i \in J\) ——生成；我们说，只有指标属于 \(J\) 的未定元出现在 \(\operatorname{Lib}_{A}(J)\) 的一个元素中。第 6 号中给出的广群代数的定义表明， \(\operatorname{Lib}_{A}(I)\) 是子代数 \(\operatorname{Lib}_{A}(J)\) 的有向族的并，其中 \(J\) 遍历 \(I\) 的有限子集的集合。对 \(\operatorname{Libas}_{A}(I)\) 与 \(\operatorname{Libasc}_{A}(I)\) 有类似的结果。
\item
  对元素 \(s\) —— \(\mathbf{M}(\mathbf{I})\) （相应地 \(\mathbf{Mo}(\mathbf{I})\) ，相应地 \(\mathbf{N}^{(1)}\) ）的——关联它的长度 \(l(s)\) ，这是一个 \(\geqslant 1\) （相应地 \(\geqslant 0\) ）的整数，使得 \(l(ss') = l(s) + l(s')\) （I，§ 7，第 1、2 与 7 号）。若 \(e_s\) 是 \(\operatorname{Lib}_{\mathbf{A}}(\mathbf{I})\) （相应地 \(\operatorname{Libas}_{\mathbf{A}}(\mathbf{I})\) ，相应地 \(\operatorname{Libasc}_{\mathbf{A}}(\mathbf{I})\) ）中对应于 \(s\) 的元素，则元素 \(x = \sum_{s} \alpha_s e_s \neq 0\) （属于 \(\operatorname{Lib}_{\mathbf{A}}(\mathbf{I})\) （相应地 \(\operatorname{Libas}_{\mathbf{A}}(\mathbf{A})\) ，相应地 \(\operatorname{Libasc}_{\mathbf{A}}(\mathbf{I})\) ））的总次数（或简称次数）是诸数 \(l(s)\) 中的最大者，其中 \(s\) 遍历（依假设非空的）满足 \(\alpha_s \neq 0\) 的元素集合。例如，若 \(i, j, k\) 是 \(\mathbf{I}\) 的三个不同元素，则元素 \((\mathrm{X}_t(\mathrm{X}_j\mathrm{X}_k))\mathrm{X}_t - (\mathrm{X}_t\mathrm{X}_j)(\mathrm{X}_k\mathrm{X}_l)\) 是一个 \(\neq 0\) 的、总次数为 4 的元素（在 \(\operatorname{Lib}_{\mathbf{A}}(\mathbf{I})\) 中）。
\end{enumerate}

